\subsection{FUNDAMENTAL WEIGHTS, DOMINANT WEIGHTS}

Assume that R is reduced. Let C be a chamber of R, and let B be the corresponding basis of R. Since R is reduced, \(B^{-} = \{\alpha^{-}\}_{\alpha \in B}\) is a basis of \(R^{-}\) . The dual basis \((\overline{\omega}_{\alpha})_{\alpha \in B}\) of \(B^{-}\) is thus a basis of the group of weights; its elements are called the fundamental weights (relative to B, or to C); if the elements of B are denoted by \((\alpha_{1}, \ldots, \alpha_{l})\) , the corresponding fundamental weights are denoted by \((\overline{\omega}_{1}, \ldots, \overline{\omega}_{l})\) .

Let \(x \in V\) . Then, \(x \in C\) if and only if \(\langle x, \alpha \rangle > 0\) for all \(\alpha \in B\) . It follows that \(C\) is the set of linear combinations of the \(\bar{\omega}_{\alpha}\) with coefficients \(\geqslant 0\) .

The entries \(n(\alpha, \beta) = \langle \alpha, \beta' \rangle\) of the Cartan matrix are, for fixed \(\alpha\) , the coordinates of \(\alpha\) with respect to the basis \((\overline{\omega}_{\beta})_{\beta \in B}\) :

\[
\alpha = \sum_ {\beta \in B} n (\alpha , \beta) \overline {{\omega}} _ {\beta}.\tag{14}
\]

The Cartan matrix is thus the transpose of the matrix of the canonical injection

\[
\mathrm{Q} (\mathrm{R}) \rightarrow \mathrm{P} (\mathrm{R})
\]

with respect to the bases B and \((\overline{\omega}_{\alpha})_{\alpha\in\mathsf{B}}\) of the Z-modules \(\mathrm{Q}(\mathrm{R})\) and \(\mathrm{P}(\mathrm{R})\) . A weight \(\overline{\omega}\) is said to be dominant if it belongs to \(\overline{C}\) , in other words if its coordinates with respect to \((\overline{\omega}_{\alpha})_{\alpha\in\mathsf{B}}\) are integers \(\geqslant0\) , or equivalently if \(g(\overline{\omega})\leqslant\overline{\omega}\) for all \(g\in\mathrm{W}(\mathrm{R})\) (no. 6, Prop. 18). Since \(\overline{C}\) is a fundamental domain for \(\mathrm{W(R)}\) (Th. 2), there exists, for any weight \(\overline{\omega}\) , a unique weight \(\overline{\omega}'\) such that \(\overline{\omega}'\) is a transform of \(\overline{\omega}\) by \(\mathrm{W(R)}\) .

We have

\[
\langle \overline {{{{\omega}}}} _ {\alpha}, \beta^ {\prime} \rangle = (\overline {{{{\omega}}}} _ {\alpha} | \frac {2 \beta}{(\beta | \beta)}) = \delta_ {\alpha \beta}
\]

for \(\alpha, \beta \in \mathbf{B}\) ( \(\delta_{\alpha\beta}\) denoting the Kronecker symbol), hence

\[
s _ {\beta} (\overline {{\omega}} _ {\alpha}) = \overline {{\omega}} _ {\alpha} - \delta_ {\alpha \beta} \beta \quad \text { and } \quad (\overline {{\omega}} _ {\alpha} | \beta) = \frac {1}{2} (\beta | \beta) \delta_ {\alpha \beta}.
\]

In other words, \(\overline{\omega}_{\alpha}\) is orthogonal to \(\beta\) for \(\beta \neq \alpha\) , and its orthogonal projection onto \(R\alpha\) is \(\frac{1}{2}\alpha\) . Since \(\overline{\omega}_{\alpha} \in \overline{C}\) , \((\overline{\omega}_{\alpha} | \overline{\omega}_{\beta}) \geqslant 0\) for \(\alpha, \beta \in B\) , i.e.~the angle \((\widehat{\overline{\omega}_{\alpha}, \overline{\omega}_{\beta}})\) is acute or a right angle. The dominant weights are the \(\overline{\omega} \in V\) such that \(2(\overline{\omega} | \alpha)/( \alpha | \alpha)\) is an integer \(\geqslant 0\) for all \(\alpha \in B\) .

PROPOSITION 28. Let B be a basis of R, B' a subset of B, V' the vector subspace of V generated by B', R' = R ∩ V' (which is a root system in V'), R' the inverse root system (which is identified with the canonical image of R' in R'), V1 the orthogonal complement of R' in V, and p the projection of V onto V' parallel to V1. Then, Q(R') = Q(R) ∩ V', P(R') = p(P(R)). The set of dominant weights of R' is the image under p of the set of dominant weights of R.

Indeed, \(\mathrm{Q}(\mathrm{R})\) is the subgroup of \(\mathrm{V}\) with basis \(\mathrm{B}\) , \(\mathrm{Q}(\mathrm{R}')\) is the subgroup of \(\mathrm{V}'\) with basis \(\mathrm{B}'\) (no. 7, Cor. 4 of Prop. 20) from which \(\mathrm{Q}(\mathrm{R}') = \mathrm{Q}(\mathrm{R}) \cap \mathrm{V}'\) is immediate. If \(\overline{\omega} \in \mathrm{P}(\mathrm{R})\) and \(\alpha \in \mathrm{R}'\) , \(\langle p(\overline{\omega}), \alpha' \rangle = \langle \overline{\omega}, \alpha' \rangle \in \mathbf{Z}\) , so \(p(\overline{\omega}) \in \mathrm{P}(\mathrm{R}')\) , and hence \(p(\mathrm{P}(\mathrm{R})) \subseteq \mathrm{P}(\mathrm{R}')\) . If \(\overline{\omega}' \in \mathrm{P}(\mathrm{R}')\) , \(\overline{\omega}'\) extends to a linear form \(\overline{\omega}\) on \(\mathrm{V}^*\) vanishing on \((\mathrm{B} - \mathrm{B}')'\) ; we have \(\langle \overline{\omega}, \alpha' \rangle \in \mathbf{Z}\) for all \(\alpha \in \mathrm{B}\) , so \(\overline{\omega} \in \mathrm{P}(\mathrm{R})\) , and \(\overline{\omega}' = p(\overline{\omega})\) ; hence \(\mathrm{P}(\mathrm{R}') \subseteq p(\mathrm{P}(\mathrm{R}))\) . Thus \(\mathrm{P}(\mathrm{R}') = p(\mathrm{P}(\mathrm{R}))\) , and the assertion about dominant weights is proved in the same way.

PROPOSITION 29. Let \(\rho\) be half the sum of the roots \(>0\) .

\begin{enumerate}
\def\labelenumi{(\roman{enumi})}
\item
  \(\rho = \sum_{\alpha \in B} \overline{\omega}_{\alpha}\) ; this is an element of C.
\item
  \(s_{\alpha}(\rho) = \rho -\alpha\) for all \(\alpha \in B\)
\item
  \((2\rho |\alpha) = (\alpha |\alpha)\) for all \(\alpha \in B\) .
\end{enumerate}

Since R is reduced, \(s_{\alpha}(R_{+}(C) - \{\alpha\}) = R_{+}(C) - \{\alpha\}\) and \(s_{\alpha}(\alpha) = -\alpha\) for \(\alpha \in B\) (no. 6, Cor. 1 of Prop. 17), so \(s_{\alpha}(2\rho) = 2\rho - 2\alpha\) . Since \(s_{\alpha}(\rho) = \rho - \langle \rho, \alpha^{-} \rangle\) , we see that

\[
\langle \rho , \alpha^ {\prime} \rangle = 1 = \left\langle \sum_ {\beta \in B} \bar {\omega} _ {\beta}, \alpha^ {\prime} \right\rangle .
\]

Hence, \(\rho = \sum_{\beta} \overline{\omega}_{\beta}\) , and consequently \(\rho \in C\) . Finally, (iii) is equivalent to \(\langle \rho, \alpha^{-} \rangle = 1\) .

COROLLARY. Let \(\sigma\) be half the sum of the elements \(>0\) of \(\mathbf{R}^{-}\) (for \(\mathbf{B}^{-}\) ). For all \(\alpha \in V\) , the sum of the coordinates of \(\alpha\) with respect to the basis \(\mathbf{B}\) is \(\langle \alpha, \sigma \rangle\) . If \(\alpha \in R\) , this sum is equal to \(\frac{1}{2} \sum_{\beta \in \mathbb{R}_{+}(C)} n(\alpha, \beta)\) .

Interchanging the roles of R and \(R^{-}\) above, we have \(\langle\alpha,\sigma\rangle=1\) for all \(\alpha\in B\) , hence the corollary.

